\clearpage

\lehead[]{\normalfont\sffamily\hspace*{-2.00cm}\textcolor{white}{\colorbox{lightblue}{\makebox[1.60cm][r]{\thechapter}}}\hspace{0.17cm}\textcolor{lightblue}{\chaptertitle}}
\rohead[]{\textcolor{lightblue}{\chaptertitle}\normalfont\sffamily\hspace*{0.17cm}\textcolor{white}{\colorbox{lightblue}{\makebox[1.60cm][l]{\thechapter}}}\hspace{-2.00cm}}
%\chead[]{}
\rehead[]{\textcolor{lightblue}{AvHG, Inf, My}}
\lohead[]{\textcolor{lightblue}{AvHG, Inf, My}}

\section{Tabellen verknüpfen: Primär- und Fremdschlüssel}

\subsection{Primärschlüssel}

Um Datensätze einer Tabelle eindeutig identifizieren zu können, braucht man
einen sogenannten \emph{Schlüssel}. Als Schlüssel kann jedes einzelne Attribut
oder -- wenn nötig -- auch eine Kombination mehrerer Attribute der Tabelle
dienen. Wichtig ist nur, dass sich der Schlüssel konzeptionell zur eindeutigen
Identifizierung eines jeden Datensatzes der Tabelle eignet. So reicht es etwa
nicht, zu sehen, dass in unserer Haustierdatenbank keine Kombination aus
Tiername und Tierart doppelt vorkommt. Denn niemand kann garantieren, dass in
Zukunft nicht doch noch ein weiteres Tier mit einer zuvor bereits
existierenden Kombination aus Tierart und Tiername hinzukommen wird. Tiername
und Tierart sind also zusammen kein geeigneter Schlüssel für die Tabelle
\myUserInput{tier}!

An einen Schlüssel wird zudem die Anforderung gestellt, dass er minimal sein
soll: Alle Attribute, die man zum Schlüssel dazuzählt, müssen auch wirklich
nötig sein, um die eindeutige Identifizierung aller Datensätze (auch
konzeptionell) sicherzustellen.

Bei einer gegebenen Tabelle können durchaus verschiedene Attribute oder
Kombinationen aus Attributen als Schlüssel infrage kommen. Den Schlüssel, den
man aus der Menge der möglichen Schlüsselkandidaten zur weiteren Verwendung
auswählt, bezeichnet man als \emph{Primärschlüssel} (engl.\ \emph{Primary Key}).

Als Primärschlüssel sind besonders gut extra für diesen Zweck erzeugte
Attribute vom Typ \lstinline|INT| geeignet, die von der Datenbank mithilfe des
Schlüsselwortes \lstinline|AUTO_INCREMENT| beim Erzeugen eines neuen
Datensatzes automatisch mit einem bisher noch nicht vergebenen Wert belegt
werden.

Überhaupt ist es empfehlenswert, als Primärschlüssel ein Attribut zu benutzen,
welches keine sonstigen Informationen trägt! Damit umgeht man so manche Falle!
Beispielsweise könnte man annehmen, dass die ISBN-Nummer ein geeigneter
Primärschlüssel für eine \myUserInput{buch}-Tabelle wäre. Oder dass die
Postleitzahl ein ebenso geeigneter Primärschlüssel für eine
\myUserInput{ort}-Tabelle wäre. Beides stellt sich bei näherer Betrachtung als
falsch heraus: Es gibt Postleitzahlen, aus denen nicht eindeutig auf einen
bestimmten Ort geschlossen werden kann, und es gibt auch ISBN-Nummern, aus
denen nicht eindeutig auf eine bestimmte Ausgabe eines Buches geschlossen werden
kann (etwa durch versehentliche Doppelvergaben von Verlagen, unveränderte
Nachdrucke oder ältere Werke ganz ohne ISBN). Durch die Verwendung eines
nicht-informationstragenden Primärschlüssels umgeht man solche Fallen. Zudem
muss man sich keine Sorgen machen, dass es in Zukunft aus inhaltlichen Gründen
notwendig werden könnte, die Werte des Primärschlüsselattributes zu ändern --
so wie es beispielsweise bei den Postleitzahlen mehrfach (zuletzt 1993)
geschehen ist.

\begin{lstlisting}
CREATE TABLE besitzer
(
  besitzer_id INT AUTO_INCREMENT,
  anrede ENUM('Herr', 'Frau', 'Firma') DEFAULT 'Herr',
  vorname VARCHAR(20),
  nachname VARCHAR(20) NOT NULL,
  PRIMARY KEY (besitzer_id)
);
\end{lstlisting}

\subsection{Fremdschlüssel}

Als \emph{Fremdschlüssel} (engl.\ \emph{Foreign Key}) bezeichnet man ein
Attribut einer Tabelle, welches auf die Werte eines Primärschlüssels einer
anderen Tabelle verweist. Genau genommen muss der Primärschlüssel noch nicht
einmal aus einer anderen Tabelle stammen: Auch Verweise auf den Primärschlüssel
der eigenen Tabelle sind möglich!

In unserer Datenbank \myUserInput{haustier} verknüpft die Tabelle
\myUserInput{beziehung} über zwei Fremdschlüssel jeweils einen Besitzer aus der
Tabelle \myUserInput{besitzer} mit einem Tier aus der Tabelle
\myUserInput{tier} (die Kombination beider Fremdschlüssel bildet hier zugleich
den zusammengesetzten Primärschlüssel der Beziehungstabelle):

\begin{lstlisting}
CREATE TABLE beziehung
(
  beziehung_besitzer_id INT NOT NULL,
  beziehung_tier_id INT NOT NULL,
  PRIMARY KEY (beziehung_besitzer_id, beziehung_tier_id),
  FOREIGN KEY (beziehung_besitzer_id) REFERENCES besitzer (besitzer_id),
  FOREIGN KEY (beziehung_tier_id) REFERENCES tier (tier_id)
);
\end{lstlisting}

Hinter \lstinline|REFERENCES| wird dabei die Tabelle und das entsprechende
Attribut benannt, auf welches sich der Fremdschlüssel bezieht. Dadurch ist es
auch sehr leicht, in einem gegebenen SQL-Skript (wie in unserer Datei
\myFile{haustier.sql}) die entsprechenden Beziehungen zwischen den Tabellen zu
erkennen.

\subsection{Referentielle Integrität}

Indem man der Datenbank mitteilt, dass es sich bei einem Attribut um einen
Fremdschlüssel handelt, versetzt man sie auch in die Lage, über die sogenannte
\emph{referentielle Integrität} der Datenbank zu wachen. Die Datenbank
verhindert dann, dass ein neuer Datensatz so angelegt wird, dass der Wert für
den Fremdschlüssel nicht durch einen bereits existierenden Wert im
Primärschlüsselfeld der referenzierten Tabelle gedeckt ist.

Im Beispiel unserer Haustierdatenbank: Nachdem wir in der Tabelle
\myUserInput{beziehung} die Attribute \linebreak \lstinline|beziehung_besitzer_id| und
\lstinline|beziehung_tier_id| als Fremdschlüssel definiert haben, wird es die
Datenbank nicht erlauben, dass wir eine Beziehung mit einer Besitzer- oder
Tier-ID eintragen, die es in den Tabellen \myUserInput{besitzer} bzw.\
\myUserInput{tier} gar nicht gibt. Ebenso würde verhindert, dass ein Datensatz
aus \myUserInput{besitzer} oder \myUserInput{tier} gelöscht oder so verändert
würde, dass in der Folge ein Verweis in der Tabelle \myUserInput{beziehung} auf
einen nicht länger existierenden Primärschlüsselwert entstünde.


\section{Tabellen verknüpfen (\texttt{JOIN}) und Unterabfragen}

\subsection{Kartesisches Produkt und \texttt{INNER JOIN}}

Gibt man hinter \lstinline|FROM| einfach zwei Tabellen ohne Verknüpfungsbedingung
an (\lstinline|SELECT * FROM besitzer, tier;|), bildet SQL das sogenannte
\emph{Kartesische Produkt}: Jede Zeile der ersten Tabelle wird mit jeder Zeile
der zweiten Tabelle kombiniert. Das Ergebnis enthält viele unsinnige
Kombinationen.

Um nur zusammengehörige Datensätze zu verbinden, verknüpft man über
\lstinline|JOIN ... ON| den Fremdschlüssel der einen Tabelle mit dem passenden
Primärschlüssel der anderen Tabelle. Da in unserer Datenbank eine
$n:m$-Beziehung vorliegt, verknüpfen wir alle drei Tabellen über die
Beziehungstabelle miteinander:

\begin{lstlisting}
SELECT besitzer.vorname, besitzer.nachname, tier.name, tier.tierart
FROM besitzer
JOIN beziehung ON besitzer.besitzer_id = beziehung.beziehung_besitzer_id
JOIN tier ON tier.tier_id = beziehung.beziehung_tier_id;
\end{lstlisting}

\subsection{\texttt{LEFT JOIN} und \texttt{RIGHT JOIN}}

Ein normaler \lstinline|JOIN| (auch \lstinline|INNER JOIN|) liefert nur Zeilen,
für die es in \emph{beiden} verknüpften Tabellen einen passenden Partner gibt.
Personen ohne Tier oder Tiere ohne Besitzer tauchen darin nicht auf.

Gegeben seien als kurzes Beispiel die beiden Tabellen \myUserInput{besitzer}
(1: Anna, 2: Andrew, 3: Annabel) und \myUserInput{fahrzeug} (Fahrrad $\rightarrow$ 2,
Auto $\rightarrow$ 4, Auto $\rightarrow$ 3):

\begin{compactitem}
  \item Ein \textbf{\lstinline|LEFT JOIN|} übernimmt \emph{alle} Zeilen der
  links stehenden Tabelle (\myUserInput{besitzer}) und füllt fehlende Partner
  rechts mit \lstinline|NULL| auf (Anna erscheint also mit \lstinline|NULL| beim
  Fahrzeug).
  \item Ein \textbf{\lstinline|RIGHT JOIN|} übernimmt umgekehrt \emph{alle}
  Zeilen der rechts stehenden Tabelle (\myUserInput{fahrzeug}) und füllt
  fehlende Partner links mit \lstinline|NULL| auf (das Auto mit Besitzer-ID $4$
  erscheint also mit \lstinline|NULL| beim Besitzer).
\end{compactitem}

\begin{lstlisting}
SELECT *
FROM besitzer LEFT JOIN fahrzeug
ON besitzer.besitzer_id = fahrzeug.fahrzeug_besitzer_id;
\end{lstlisting}

\subsection{Verknüpfung einer Tabelle mit sich selbst (Self-Join)}

Gegeben sei eine Tabelle \myUserInput{angestellter} mit den Spalten
\myUserInput{name} und \myUserInput{job} (Jan Bauer: Buchhalter; Talina Wendel:
Programmierer; Noah Richter: Programmierer; Andrea Faust: Redakteur). Um
herauszufinden, wer denselben Job hat wie Talina Wendel, verwendet man dieselbe
Tabelle zweimal und vergibt zur Unterscheidung mithilfe von \lstinline|AS|
sogenannte \emph{Aliase}:

\begin{lstlisting}
SELECT kollege.name
FROM angestellter AS talina
JOIN angestellter AS kollege ON kollege.job = talina.job
WHERE talina.name = 'Talina Wendel';
\end{lstlisting}

\noindent \textbf{Denkaufgabe:} Welche Ausgabe erzeugt diese Abfrage?

\textbf{Hinweis:} Self-Joins, also die Dopplung von Tabellen durch Alias-Namen mit 
\lstinline|AS|, werden schnell kompliziert - insbesondere dann, wenn nicht nur eine, 
sondern mehrere Tabellen auf diese Weise gedoppelt werden müssen. In vielen Fällen 
(immer dann, wenn nur nach übereinstimmenden Werten gefiltert wird) lässt sich ein 
Self-Join auch leichter als Unterabfrage mit \lstinline|=| oder \lstinline|IN| 
schreiben!

\subsection{Unterabfragen (Subqueries) mit \texttt{=} und \texttt{IN}}

Häufig möchte man Daten aus einer Tabelle mithilfe eines Wertes filtern, den man
selbst erst durch eine Abfrage ermitteln muss. Dazu kann man eine
\lstinline|SELECT|-Abfrage direkt innerhalb einer anderen Abfrage verwenden
(sogenannte \emph{Unterabfrage} oder \emph{Subquery}).

\subsubsection{Einfacher Vergleich mit \texttt{=}}

Liefert die innere Abfrage garantiert \textbf{genau einen einzelnen Wert} zurück,
kann man den normalen Gleichheitsoperator \lstinline|=| verwenden.
Beispielsweise lassen sich so alle Tiere ermitteln, die am selben Tag geboren
wurden wie das Tier \myUserInput{Maja}:

\begin{lstlisting}
SELECT name
FROM tier
WHERE geburtstag = (
    SELECT geburtstag
    FROM tier
    WHERE name = 'Maja'
)
AND name != 'Maja';
\end{lstlisting}

\subsubsection{Mehrere Rückgabewerte mit \texttt{IN}}

Kann die innere Abfrage \textbf{mehrere Werte} zurückliefern (z.\,B.\ weil ein Tier
mehrere Besitzer hat oder es mehrere Personen mit demselben Job gibt), würde
\lstinline|=| zu einem Fehler führen. Stattdessen verwendet man den Operator
\lstinline|IN|:

\begin{lstlisting}
SELECT name
FROM angestellter
WHERE job IN (
    SELECT job
    FROM angestellter
    WHERE name = 'Talina Wendel'
);
\end{lstlisting}

Schritt für Schritt passiert hier Folgendes: Die innere Abfrage ermittelt
zunächst den Job (oder die Jobs) von Talina Wendel. Die äußere Abfrage gibt
anschließend alle Angestellten aus, deren Job in dieser Ergebnisliste enthalten
ist.
